% ===========================================================================
%  Chapter 5 -- Discussion and Conclusions
%  Disease-Adaptive Heterogeneous Graph Learning with Anatomically Aligned
%  Multi-Sequence MRI Representations
%
%  DRAFT 1 -- written 2026-08-27 against the completed three-seed campaign, the
%  controlled input ablation, the attribution analysis and the ROI quality
%  control. Figures tagged [3-SEED] refresh when the seven-seed campaign lands;
%  the argument of this chapter does not depend on them, because every
%  comparison it leans on either holds on all seeds or is reported as null.
%
%  Compiles standalone:   pdflatex chapter5 ; biber chapter5 ; pdflatex chapter5
%  To fold into the thesis, % ===========================================================================
%  Chapter 5 -- Discussion and Conclusions
%  Disease-Adaptive Heterogeneous Graph Learning with Anatomically Aligned
%  Multi-Sequence MRI Representations
%
%  DRAFT 1 -- written 2026-08-27 against the completed three-seed campaign, the
%  controlled input ablation, the attribution analysis and the ROI quality
%  control. Figures tagged [3-SEED] refresh when the seven-seed campaign lands;
%  the argument of this chapter does not depend on them, because every
%  comparison it leans on either holds on all seeds or is reported as null.
%
%  Compiles standalone:   pdflatex chapter5 ; biber chapter5 ; pdflatex chapter5
%  To fold into the thesis, % ===========================================================================
%  Chapter 5 -- Discussion and Conclusions
%  Disease-Adaptive Heterogeneous Graph Learning with Anatomically Aligned
%  Multi-Sequence MRI Representations
%
%  DRAFT 1 -- written 2026-08-27 against the completed three-seed campaign, the
%  controlled input ablation, the attribution analysis and the ROI quality
%  control. Figures tagged [3-SEED] refresh when the seven-seed campaign lands;
%  the argument of this chapter does not depend on them, because every
%  comparison it leans on either holds on all seeds or is reported as null.
%
%  Compiles standalone:   pdflatex chapter5 ; biber chapter5 ; pdflatex chapter5
%  To fold into the thesis, % ===========================================================================
%  Chapter 5 -- Discussion and Conclusions
%  Disease-Adaptive Heterogeneous Graph Learning with Anatomically Aligned
%  Multi-Sequence MRI Representations
%
%  DRAFT 1 -- written 2026-08-27 against the completed three-seed campaign, the
%  controlled input ablation, the attribution analysis and the ROI quality
%  control. Figures tagged [3-SEED] refresh when the seven-seed campaign lands;
%  the argument of this chapter does not depend on them, because every
%  comparison it leans on either holds on all seeds or is reported as null.
%
%  Compiles standalone:   pdflatex chapter5 ; biber chapter5 ; pdflatex chapter5
%  To fold into the thesis, \input{chapter5}; docmute skips this preamble.
% ===========================================================================

\documentclass[12pt,a4paper]{report}

\usepackage{lmodern}
\usepackage[T1]{fontenc}
\usepackage[utf8]{inputenc}
\usepackage[english]{babel}
\usepackage{csquotes}
\usepackage{microtype}
\usepackage[margin=2.5cm]{geometry}
\usepackage{setspace}
\onehalfspacing
\usepackage{amsmath,amssymb}
\usepackage{booktabs}
\usepackage{array}
\usepackage{graphicx}
\usepackage[hidelinks]{hyperref}

\begin{document}

\chapter{Discussion and Conclusions}
\label{ch:discussion}

% ===========================================================================
\section{What This Study Found}
\label{sec:disc-summary}

The thesis asked whether explicit anatomical structure, imposed on a
multi-sequence grading model, improves automated classification of lumbar
degenerative disease. The answer is that it largely does not, and the more
useful part of the answer is \emph{why}.

The complete system improves on a single-sequence baseline by $+0.0172$ QWK
(95\% CI $[+0.0064, +0.0285]$), reproducibly across seeds and after correction
for multiple comparisons. But the components credited with that improvement in
the original design are not the components responsible for it. Anatomically
aligned self-supervision produced no measurable benefit on any of three axes.
Disease-conditioned routing learned the clinically expected sequence weights
and produced no benefit from having learned them. Relational typing helped;
the anatomical identity of the relations did not.

Two of the three proposed mechanisms are therefore rejected, one is
half-supported in a narrower form than proposed, and the largest single
contributor to final performance is an objective function that
Chapter~\ref{ch:introduction} explicitly declined to claim as original.

This chapter argues that these results are informative rather than merely
disappointing, offers two mechanisms for them, situates them against the
literature of Chapter~\ref{ch:litreview}, and states plainly what the study
cannot conclude.

% ===========================================================================
\section{Why the Anatomical Priors Did Not Pay}
\label{sec:disc-why}

A null result without a mechanism invites the reader to supply their own, and
the most available one---that the implementation was faulty---is the least
interesting and the most damaging. Two mechanisms are supported by evidence
gathered specifically to test them.

\subsection{The Encoder Already Localises}
\label{sec:disc-localise}

The structural priors of RQ1--RQ3 share a premise: that anatomical context is
needed to find and interpret the graded structure, and that supplying it
explicitly will help. The attribution analysis of
Section~\ref{sec:results-attribution} tests that premise directly and finds it
largely unmet.

Every trained configuration concentrates between $2.3$ and $2.9$ times the
chance share of its attention on the annotated lesion. Decisively, E0---a plain
single-sequence convolutional network with no routing, no self-supervision and
no graph---already reaches $0.490$ against an untrained floor of $0.204$ and a
uniform expectation of $0.197$. The localisation the priors were designed to
supply is largely accomplished before any of them is added.

This is consistent with Section~\ref{sec:lr-localize-first}, which
assembles the evidence that localisation precedes classification. The present
result sharpens that literature rather than contradicting it: localisation does
precede classification, and on a benchmark that supplies the localisation as an
annotated coordinate, a convolutional encoder given a tight physical crop
performs the residual localisation unaided. Structure imposed above that
encoder is competing for a job already done.

Section~\ref{sec:intro-conceptual} anticipated this possibility in advance,
stating that \enquote{if local evidence alone is sufficient, more complex
mechanisms should not be retained}. The finding is that local evidence is very
nearly sufficient.

\subsection{The Reference Standard Constrains the Headroom}
\label{sec:disc-ceiling}

Section~\ref{sec:lr-label-ceiling} argues that reader reliability is a hard
constraint on interpretation rather than a peripheral caveat. Lurie et al.\
report inter-reader $\kappa$ of $0.73$ for central canal stenosis, $0.58$ for
foraminal stenosis and $0.49$ for subarticular stenosis
\cite{lurie2008reliability}. A model trained on one reader's labels cannot be
expected to exceed the agreement those readers achieve with each other.

The comparison must be made carefully. Quadratic weighted kappa credits
near-misses that unweighted $\kappa$ penalises fully, so the QWK values reported
here are not directly comparable with Lurie's figures, and the two are measured
on different cohorts under different protocols. The point is not that this
system outperforms radiologists; nothing in this study supports that claim.

The point is about headroom. If the achievable band is bounded by a reference
standard on which experts disagree at $\kappa \approx 0.5$ for the lateral
compartments, then the margin available to any architectural refinement is
narrow, and effects of the size proposed here---$0.005$ QWK---sit far inside the
noise of the labels themselves. Chapter~\ref{ch:litreview} warns that
\enquote{a reported accuracy approaching unity on such a target should prompt
scrutiny of the evaluation protocol rather than admiration}; the corollary is
that failure to improve on such a target by half a percentage point warrants
less alarm than it first appears to.

\subsection{A Test of That Explanation, and Its Limits}
\label{sec:disc-ordering}

Section~\ref{sec:lr-reliability} makes a checkable prediction: agreement is
\enquote{systematically worse for the lateral compartments than for the central
canal---the same ordering that later appears in automated model performance}.
If model difficulty tracks reader difficulty, the label-ceiling explanation
gains support.

\begin{table}[htbp]
\centering
\caption{Per-condition performance against reported inter-reader agreement.
Reader $\kappa$ from Lurie et al.; model QWK averaged over three seeds. The two
metrics are not directly comparable in absolute terms; only the ordering is
being examined.}
\label{tab:disc-ordering}
\begin{tabular}{lcccc}
\toprule
Condition & Reader $\kappa$ & E0 & E2 & E4 \\
\midrule
Central canal        & 0.73 & 0.799 & 0.786 & 0.805 \\
Left foraminal       & 0.58 & 0.659 & 0.623 & 0.644 \\
Right foraminal      & 0.58 & 0.685 & 0.672 & 0.671 \\
Left subarticular    & 0.49 & 0.758 & 0.749 & 0.767 \\
Right subarticular   & 0.49 & 0.712 & 0.731 & 0.736 \\
\midrule
Central $-$ lateral mean & $+0.20$ & $+0.095$ & $+0.092$ & $+0.100$ \\
\bottomrule
\end{tabular}
\end{table}

The coarse prediction holds. Central canal stenosis is the easiest condition for
the model, as it is for readers, and the central-versus-lateral gap runs in the
same direction for both, though smaller in magnitude for the model.

The fine ordering does not hold, and the discrepancy is informative. Readers find
\emph{subarticular} stenosis hardest at $\kappa = 0.49$; the model finds
\emph{foraminal} stenosis hardest, while subarticular is its second-best
condition. Model difficulty is therefore not a simple reflection of label noise.

A representational explanation fits the discrepancy. Foraminal targets are
graded on sagittal T1, where left--right is a \emph{through-plane} axis: the two
sides are different slices rather than different in-plane positions. Subarticular
targets are graded on axial T2, where left--right is in-plane and the lateral
recesses are directly visible. The 2.5D representation used here---three adjacent
slices treated as channels---captures in-plane structure well and through-plane
structure poorly, which is exactly the axis on which foraminal grading depends.
If correct, this locates the residual difficulty in the input representation
rather than in the absence of anatomical priors, and it predicts that a volumetric
or cross-plane-aware representation would help foraminal grading specifically.

That prediction is not tested here. It is offered as the most economical account
of a dissociation the study did not set out to produce, and it is the clearest
direction for subsequent work.

Two cautions apply to this section. Only five conditions exist, two pairs of
which share a reported $\kappa$, so rank correlation between the two orderings
has essentially no statistical power and none is claimed. And Lurie's cohort is
not LumbarDISC; the comparison is indicative rather than a matched analysis.

% ===========================================================================
\section{The Components Individually}
\label{sec:disc-components}

\subsection{Anatomical Self-Supervision (RQ2)}

The pretext task was learnable and was learned: validation InfoNCE reached
$1.5162$ against a chance value of $3.4657$, so the encoders demonstrably
learned to match corresponding anatomy across sequences. The representation
that resulted conferred no measurable downstream benefit on accuracy, on
robustness to a withheld sequence, or on where the model attends.

The natural objection is that the correspondence itself was wrong---that the
pretext task matched misaligned patches and the null is an artefact. That
objection is answered. Section~\ref{sec:results-roiqc} verifies the DICOM
correspondence externally, comparing two independent human annotations through
the transform across 9{,}542 projections, of which 94\% on the pretraining
partition land within one slice of an independently annotated target at the same
level. ACSSL was given correctly aligned data and still conferred nothing.

The result is best read alongside Section~\ref{sec:disc-localise}. Cross-sequence
self-supervision is designed to teach an encoder that one anatomy appears under
several acquisitions. If the supervised task already drives the encoder to
localise that anatomy, the pretext task is teaching something the downstream
objective supplies for free.

\subsection{Disease-Conditioned Routing (RQ3)}

The routing result is the study's clearest illustration of why a mechanism check
requires a control. The router reproduces the clinically expected allocation in
all fifteen runs in which a gate is present, and withholding the
high-weight sequence for a condition is severely damaging---between $0.06$ and
$0.35$ QWK, an order of magnitude larger than any difference in the ablation
ladder. Presented alone, that is a compelling causal story.

It does not survive its control. E1, which has no router at all, exhibits the
same dependence with drops as large or larger. The pattern belongs to the
dataset: each condition is annotated on one particular sequence, so withholding
that sequence removes the only crop centred on the pathology. Any model suffers,
routed or not. Isolating the router's own contribution to sequence reliance
gives $-0.0010 \pm 0.0382$ across seeds, indistinguishable from zero.

Chapter~\ref{ch:methodology} anticipated precisely this trap, stating that
\enquote{the gate weights are model allocations, not causal estimates of
sequence importance}, and requiring the intervention that produced the control.
The methodology's caution was correct and load-bearing.

What robustness the system does exhibit under missing sequences comes from
modality dropout, which reduces reliance on the annotated sequence by
$-0.0588$ on every seed. That is a standard regulariser, disclaimed in
Section~\ref{sec:intro-contrib-supporting} as not original, and it outperforms
the mechanism designed for the purpose.

\subsection{Relational Structure (RQ1)}

RQ1 divides, and the division is the more interesting result. Typed
heterogeneous edges improve on a homogeneous graph by $+0.0123$ QWK on every
seed. Anatomically correct topology does not improve on a degree-preserving
random shuffle: $+0.0051$, two seeds of three, well inside the corrected
threshold. The control preserves node count, edge count and degree, so the two
models are matched in capacity and differ only in \emph{which} targets are
connected to \emph{which}.

The conclusion is therefore narrower and sharper than H3 proposed. What the
graph supplies is a structured parameterisation of interaction between
targets---additional capacity, organised by relation type---rather than the
specific adjacency of levels, compartments and bilateral pairs that motivated
it. Relational typing carries information; anatomical adjacency does not.

Section~\ref{sec:lr-graphs} notes that relational and graph-based modelling of
lumbar targets remains sparsely explored. The present result suggests that where
it is explored, the shuffled-topology control should be mandatory, because a
typed graph will outperform a homogeneous one for reasons that have nothing to
do with the anatomy its edges are named after.

\subsection{Ordinal and Cost-Sensitive Objectives (RQ5)}

RQ5 is the study's most reliable positive result and it contradicts the strongest
prior evidence. Section~\ref{sec:lr-ordinal-losses} reports that Niemeyer et al.\
tested soft kappa loss, ordinal cross-entropy and regression losses against
standard cross-entropy for Pfirrmann grading across 7{,}948 discs with extensive
hyperparameter optimisation, and that none improved classification performance
\cite{niemeyer2021a}.

Here the ordinal and cost-sensitive head improves QWK by $+0.0099$ on every seed,
with the smallest between-seed variance of any comparison in the study
($0.0017$), and raises recall on Severe targets from $59.6\%$ to $65.0\%$
relative to the same architecture under categorical cross-entropy.

The difference from Niemeyer et al.\ is most plausibly the cost matrix rather
than the ordinal structure. Their comparison isolated ordinality; the objective
tested here additionally encodes clinical asymmetry, penalising
Severe$\rightarrow$Normal errors more heavily than the reverse. The gain appears
where that asymmetry acts. This reading is consistent with
Section~\ref{sec:lr-error-asymmetry}, which argues that the clinical cost of an
error depends on its direction while standard objectives optimise a symmetric
one.

It is worth stating plainly that the supporting component outperformed the novel
ones, and that this is a result about where effort is best spent rather than an
embarrassment.

% ===========================================================================
\section{Methodological Implications}
\label{sec:disc-method}

Three analyses in this study produced positive conclusions that their controls
subsequently overturned. All three erred in the direction of the hypothesis
under test. They are reported in Section~\ref{sec:results-threats} and are
summarised here because the pattern is more instructive than any individual
case.

\emph{Per-seed confidence intervals understate variance.} A paired bootstrap
that resamples patients with the trained model held fixed excludes training
stochasticity. On this problem that omission is decisive: one comparison
returned $+0.0270$ with $p = 0.000$ on one seed and $-0.0234$ with $p = 0.000$
on another. Reported per seed, any claim in this study could have been supported
by choosing a seed.

\emph{A mechanism check without a control reads as proof.} The routing
intervention produced 5/5 agreement with effects an order of magnitude above any
ladder difference, and a router-free model reproduced it exactly.

\emph{An instrument can penalise the hypothesis it is measuring.} The first
attribution probe hooked the first encoder unconditionally, which for
multi-sequence configurations reads the sagittal T1 encoder even on targets
graded from axial T2. It reported that no configuration improved on the
baseline; hooking the encoder that actually read each target reverses the sign.

The general lesson is that in a study of this kind the controls are not
supporting apparatus but the primary instrument. A degree-preserving shuffle, a
router-free arm, an architecture-matched untrained floor and a label-shuffled
pipeline control are what separate a mechanism from a coincidence, and each of
them changed a conclusion here.

% ===========================================================================
\section{Limitations}
\label{sec:disc-limits}

\emph{Statistical power.} The study is underpowered for single-component effects
of the size observed. Table~\ref{tab:power} converts each effect into the number
of training seeds required to detect it: one for the full system and for the
ordinal head, ten for typed edges, but 216 for cross-sequence self-supervision
and 3{,}467 for routing. The last is roughly three years of continuous
computation on the hardware used. Those two components are not merely unproven
here; effects of the size observed are not establishable at any plausible scale,
and would carry no clinical meaning if they were.

\emph{External validity is untested.} RQ4 was not executed. The complete system
under institutional transfer, specified as rung E8, was not run. Every result in
this thesis is internal to LumbarDISC, and no claim of generalisation to another
institution, scanner population or reference standard is made.

\emph{Single benchmark, single reference standard.} LumbarDISC labels derive
from a reading process whose inter-reader agreement is moderate for the lateral
compartments. Conclusions about which mechanisms help are conditional on that
reference standard.

\emph{Attribution is not measured on the final system.} The Grad-CAM analysis
covers E0--E4; the graph rungs route through a different batch layout. The
mechanism offered in Section~\ref{sec:disc-localise} is therefore established on
the encoder rather than on E7 in full.

\emph{No clinical outcome is measured.} The study measures agreement with a
retrospective reference standard. It does not measure diagnostic accuracy
against a surgical or clinical ground truth, does not measure effect on
reporting time, and does not support any claim of patient benefit.

\emph{The quality-control reader pass is outstanding.} The ROI inspection
sheets and the checklist of Section~\ref{sec:results-roiqc} have been generated;
the reader adjudication that converts them into the promised
inclusion/exclusion table has not been completed.

% ===========================================================================
\section{Implications and Future Work}
\label{sec:disc-future}

\emph{For anatomy-aware model design.} The most transferable finding is that a
convolutional encoder given a tight, physically-defined crop already performs
most of the localisation that anatomical priors are introduced to supply. Before
adding structure above such an encoder, the marginal room for that structure
should be measured---an attribution floor against an untrained network costs
nothing and would have predicted these results in advance.

\emph{For relational modelling.} A degree-preserving shuffle should be standard
practice wherever a graph's topology is claimed to be meaningful. Without it,
capacity and anatomy are confounded, and this study finds that the capacity
explains the gain.

\emph{For the representation.} The dissociation of
Section~\ref{sec:disc-ordering}---readers find subarticular hardest, the model
finds foraminal hardest---points to the 2.5D representation as a specific,
testable limitation on through-plane structure. A volumetric or genuinely
cross-plane representation is the most direct follow-up this study suggests.

\emph{For transfer.} The infrastructure for RQ4 exists: the cohort is
reconciled, reports are parsed, and the geometric machinery is verified. What
remains is reader adjudication of conflicting reference labels, a correct
de-identification procedure, and---the largest gap---localisation for a cohort
that carries no annotated coordinates. Automated disc localisation is the
enabling step, and it is a separate piece of work rather than an extension of
this one.

% ===========================================================================
\section{Conclusion}
\label{sec:disc-conclusion}

This thesis set out to determine whether explicit anatomical structure improves
automated grading of lumbar degenerative disease. It determined that, at this
data scale and under this reference standard, it largely does not, and it
established why: the convolutional encoder already localises the graded
structure, and the reference standard leaves narrow headroom above it.

The system that was built works, reproducibly and by a margin that survives
correction. The components that make it work are relational typing and a
clinically motivated objective, not anatomical correspondence or
disease-conditioned routing. Two proposed mechanisms are rejected on evidence
gathered specifically to test them, a third is supported in a narrower and more
precise form than proposed, and the study's own controls overturned three
conclusions that would otherwise have been reported as findings.

Chapter~\ref{ch:introduction} framed the aim as \enquote{a determination rather
than a promise of superiority}. That framing was prescient, and this thesis
should be read as having discharged it.

\end{document}
; docmute skips this preamble.
% ===========================================================================

\documentclass[12pt,a4paper]{report}

\usepackage{lmodern}
\usepackage[T1]{fontenc}
\usepackage[utf8]{inputenc}
\usepackage[english]{babel}
\usepackage{csquotes}
\usepackage{microtype}
\usepackage[margin=2.5cm]{geometry}
\usepackage{setspace}
\onehalfspacing
\usepackage{amsmath,amssymb}
\usepackage{booktabs}
\usepackage{array}
\usepackage{graphicx}
\usepackage[hidelinks]{hyperref}

\begin{document}

\chapter{Discussion and Conclusions}
\label{ch:discussion}

% ===========================================================================
\section{What This Study Found}
\label{sec:disc-summary}

The thesis asked whether explicit anatomical structure, imposed on a
multi-sequence grading model, improves automated classification of lumbar
degenerative disease. The answer is that it largely does not, and the more
useful part of the answer is \emph{why}.

The complete system improves on a single-sequence baseline by $+0.0172$ QWK
(95\% CI $[+0.0064, +0.0285]$), reproducibly across seeds and after correction
for multiple comparisons. But the components credited with that improvement in
the original design are not the components responsible for it. Anatomically
aligned self-supervision produced no measurable benefit on any of three axes.
Disease-conditioned routing learned the clinically expected sequence weights
and produced no benefit from having learned them. Relational typing helped;
the anatomical identity of the relations did not.

Two of the three proposed mechanisms are therefore rejected, one is
half-supported in a narrower form than proposed, and the largest single
contributor to final performance is an objective function that
Chapter~\ref{ch:introduction} explicitly declined to claim as original.

This chapter argues that these results are informative rather than merely
disappointing, offers two mechanisms for them, situates them against the
literature of Chapter~\ref{ch:litreview}, and states plainly what the study
cannot conclude.

% ===========================================================================
\section{Why the Anatomical Priors Did Not Pay}
\label{sec:disc-why}

A null result without a mechanism invites the reader to supply their own, and
the most available one---that the implementation was faulty---is the least
interesting and the most damaging. Two mechanisms are supported by evidence
gathered specifically to test them.

\subsection{The Encoder Already Localises}
\label{sec:disc-localise}

The structural priors of RQ1--RQ3 share a premise: that anatomical context is
needed to find and interpret the graded structure, and that supplying it
explicitly will help. The attribution analysis of
Section~\ref{sec:results-attribution} tests that premise directly and finds it
largely unmet.

Every trained configuration concentrates between $2.3$ and $2.9$ times the
chance share of its attention on the annotated lesion. Decisively, E0---a plain
single-sequence convolutional network with no routing, no self-supervision and
no graph---already reaches $0.490$ against an untrained floor of $0.204$ and a
uniform expectation of $0.197$. The localisation the priors were designed to
supply is largely accomplished before any of them is added.

This is consistent with Section~\ref{sec:lr-localize-first}, which
assembles the evidence that localisation precedes classification. The present
result sharpens that literature rather than contradicting it: localisation does
precede classification, and on a benchmark that supplies the localisation as an
annotated coordinate, a convolutional encoder given a tight physical crop
performs the residual localisation unaided. Structure imposed above that
encoder is competing for a job already done.

Section~\ref{sec:intro-conceptual} anticipated this possibility in advance,
stating that \enquote{if local evidence alone is sufficient, more complex
mechanisms should not be retained}. The finding is that local evidence is very
nearly sufficient.

\subsection{The Reference Standard Constrains the Headroom}
\label{sec:disc-ceiling}

Section~\ref{sec:lr-label-ceiling} argues that reader reliability is a hard
constraint on interpretation rather than a peripheral caveat. Lurie et al.\
report inter-reader $\kappa$ of $0.73$ for central canal stenosis, $0.58$ for
foraminal stenosis and $0.49$ for subarticular stenosis
\cite{lurie2008reliability}. A model trained on one reader's labels cannot be
expected to exceed the agreement those readers achieve with each other.

The comparison must be made carefully. Quadratic weighted kappa credits
near-misses that unweighted $\kappa$ penalises fully, so the QWK values reported
here are not directly comparable with Lurie's figures, and the two are measured
on different cohorts under different protocols. The point is not that this
system outperforms radiologists; nothing in this study supports that claim.

The point is about headroom. If the achievable band is bounded by a reference
standard on which experts disagree at $\kappa \approx 0.5$ for the lateral
compartments, then the margin available to any architectural refinement is
narrow, and effects of the size proposed here---$0.005$ QWK---sit far inside the
noise of the labels themselves. Chapter~\ref{ch:litreview} warns that
\enquote{a reported accuracy approaching unity on such a target should prompt
scrutiny of the evaluation protocol rather than admiration}; the corollary is
that failure to improve on such a target by half a percentage point warrants
less alarm than it first appears to.

\subsection{A Test of That Explanation, and Its Limits}
\label{sec:disc-ordering}

Section~\ref{sec:lr-reliability} makes a checkable prediction: agreement is
\enquote{systematically worse for the lateral compartments than for the central
canal---the same ordering that later appears in automated model performance}.
If model difficulty tracks reader difficulty, the label-ceiling explanation
gains support.

\begin{table}[htbp]
\centering
\caption{Per-condition performance against reported inter-reader agreement.
Reader $\kappa$ from Lurie et al.; model QWK averaged over three seeds. The two
metrics are not directly comparable in absolute terms; only the ordering is
being examined.}
\label{tab:disc-ordering}
\begin{tabular}{lcccc}
\toprule
Condition & Reader $\kappa$ & E0 & E2 & E4 \\
\midrule
Central canal        & 0.73 & 0.799 & 0.786 & 0.805 \\
Left foraminal       & 0.58 & 0.659 & 0.623 & 0.644 \\
Right foraminal      & 0.58 & 0.685 & 0.672 & 0.671 \\
Left subarticular    & 0.49 & 0.758 & 0.749 & 0.767 \\
Right subarticular   & 0.49 & 0.712 & 0.731 & 0.736 \\
\midrule
Central $-$ lateral mean & $+0.20$ & $+0.095$ & $+0.092$ & $+0.100$ \\
\bottomrule
\end{tabular}
\end{table}

The coarse prediction holds. Central canal stenosis is the easiest condition for
the model, as it is for readers, and the central-versus-lateral gap runs in the
same direction for both, though smaller in magnitude for the model.

The fine ordering does not hold, and the discrepancy is informative. Readers find
\emph{subarticular} stenosis hardest at $\kappa = 0.49$; the model finds
\emph{foraminal} stenosis hardest, while subarticular is its second-best
condition. Model difficulty is therefore not a simple reflection of label noise.

A representational explanation fits the discrepancy. Foraminal targets are
graded on sagittal T1, where left--right is a \emph{through-plane} axis: the two
sides are different slices rather than different in-plane positions. Subarticular
targets are graded on axial T2, where left--right is in-plane and the lateral
recesses are directly visible. The 2.5D representation used here---three adjacent
slices treated as channels---captures in-plane structure well and through-plane
structure poorly, which is exactly the axis on which foraminal grading depends.
If correct, this locates the residual difficulty in the input representation
rather than in the absence of anatomical priors, and it predicts that a volumetric
or cross-plane-aware representation would help foraminal grading specifically.

That prediction is not tested here. It is offered as the most economical account
of a dissociation the study did not set out to produce, and it is the clearest
direction for subsequent work.

Two cautions apply to this section. Only five conditions exist, two pairs of
which share a reported $\kappa$, so rank correlation between the two orderings
has essentially no statistical power and none is claimed. And Lurie's cohort is
not LumbarDISC; the comparison is indicative rather than a matched analysis.

% ===========================================================================
\section{The Components Individually}
\label{sec:disc-components}

\subsection{Anatomical Self-Supervision (RQ2)}

The pretext task was learnable and was learned: validation InfoNCE reached
$1.5162$ against a chance value of $3.4657$, so the encoders demonstrably
learned to match corresponding anatomy across sequences. The representation
that resulted conferred no measurable downstream benefit on accuracy, on
robustness to a withheld sequence, or on where the model attends.

The natural objection is that the correspondence itself was wrong---that the
pretext task matched misaligned patches and the null is an artefact. That
objection is answered. Section~\ref{sec:results-roiqc} verifies the DICOM
correspondence externally, comparing two independent human annotations through
the transform across 9{,}542 projections, of which 94\% on the pretraining
partition land within one slice of an independently annotated target at the same
level. ACSSL was given correctly aligned data and still conferred nothing.

The result is best read alongside Section~\ref{sec:disc-localise}. Cross-sequence
self-supervision is designed to teach an encoder that one anatomy appears under
several acquisitions. If the supervised task already drives the encoder to
localise that anatomy, the pretext task is teaching something the downstream
objective supplies for free.

\subsection{Disease-Conditioned Routing (RQ3)}

The routing result is the study's clearest illustration of why a mechanism check
requires a control. The router reproduces the clinically expected allocation in
all fifteen runs in which a gate is present, and withholding the
high-weight sequence for a condition is severely damaging---between $0.06$ and
$0.35$ QWK, an order of magnitude larger than any difference in the ablation
ladder. Presented alone, that is a compelling causal story.

It does not survive its control. E1, which has no router at all, exhibits the
same dependence with drops as large or larger. The pattern belongs to the
dataset: each condition is annotated on one particular sequence, so withholding
that sequence removes the only crop centred on the pathology. Any model suffers,
routed or not. Isolating the router's own contribution to sequence reliance
gives $-0.0010 \pm 0.0382$ across seeds, indistinguishable from zero.

Chapter~\ref{ch:methodology} anticipated precisely this trap, stating that
\enquote{the gate weights are model allocations, not causal estimates of
sequence importance}, and requiring the intervention that produced the control.
The methodology's caution was correct and load-bearing.

What robustness the system does exhibit under missing sequences comes from
modality dropout, which reduces reliance on the annotated sequence by
$-0.0588$ on every seed. That is a standard regulariser, disclaimed in
Section~\ref{sec:intro-contrib-supporting} as not original, and it outperforms
the mechanism designed for the purpose.

\subsection{Relational Structure (RQ1)}

RQ1 divides, and the division is the more interesting result. Typed
heterogeneous edges improve on a homogeneous graph by $+0.0123$ QWK on every
seed. Anatomically correct topology does not improve on a degree-preserving
random shuffle: $+0.0051$, two seeds of three, well inside the corrected
threshold. The control preserves node count, edge count and degree, so the two
models are matched in capacity and differ only in \emph{which} targets are
connected to \emph{which}.

The conclusion is therefore narrower and sharper than H3 proposed. What the
graph supplies is a structured parameterisation of interaction between
targets---additional capacity, organised by relation type---rather than the
specific adjacency of levels, compartments and bilateral pairs that motivated
it. Relational typing carries information; anatomical adjacency does not.

Section~\ref{sec:lr-graphs} notes that relational and graph-based modelling of
lumbar targets remains sparsely explored. The present result suggests that where
it is explored, the shuffled-topology control should be mandatory, because a
typed graph will outperform a homogeneous one for reasons that have nothing to
do with the anatomy its edges are named after.

\subsection{Ordinal and Cost-Sensitive Objectives (RQ5)}

RQ5 is the study's most reliable positive result and it contradicts the strongest
prior evidence. Section~\ref{sec:lr-ordinal-losses} reports that Niemeyer et al.\
tested soft kappa loss, ordinal cross-entropy and regression losses against
standard cross-entropy for Pfirrmann grading across 7{,}948 discs with extensive
hyperparameter optimisation, and that none improved classification performance
\cite{niemeyer2021a}.

Here the ordinal and cost-sensitive head improves QWK by $+0.0099$ on every seed,
with the smallest between-seed variance of any comparison in the study
($0.0017$), and raises recall on Severe targets from $59.6\%$ to $65.0\%$
relative to the same architecture under categorical cross-entropy.

The difference from Niemeyer et al.\ is most plausibly the cost matrix rather
than the ordinal structure. Their comparison isolated ordinality; the objective
tested here additionally encodes clinical asymmetry, penalising
Severe$\rightarrow$Normal errors more heavily than the reverse. The gain appears
where that asymmetry acts. This reading is consistent with
Section~\ref{sec:lr-error-asymmetry}, which argues that the clinical cost of an
error depends on its direction while standard objectives optimise a symmetric
one.

It is worth stating plainly that the supporting component outperformed the novel
ones, and that this is a result about where effort is best spent rather than an
embarrassment.

% ===========================================================================
\section{Methodological Implications}
\label{sec:disc-method}

Three analyses in this study produced positive conclusions that their controls
subsequently overturned. All three erred in the direction of the hypothesis
under test. They are reported in Section~\ref{sec:results-threats} and are
summarised here because the pattern is more instructive than any individual
case.

\emph{Per-seed confidence intervals understate variance.} A paired bootstrap
that resamples patients with the trained model held fixed excludes training
stochasticity. On this problem that omission is decisive: one comparison
returned $+0.0270$ with $p = 0.000$ on one seed and $-0.0234$ with $p = 0.000$
on another. Reported per seed, any claim in this study could have been supported
by choosing a seed.

\emph{A mechanism check without a control reads as proof.} The routing
intervention produced 5/5 agreement with effects an order of magnitude above any
ladder difference, and a router-free model reproduced it exactly.

\emph{An instrument can penalise the hypothesis it is measuring.} The first
attribution probe hooked the first encoder unconditionally, which for
multi-sequence configurations reads the sagittal T1 encoder even on targets
graded from axial T2. It reported that no configuration improved on the
baseline; hooking the encoder that actually read each target reverses the sign.

The general lesson is that in a study of this kind the controls are not
supporting apparatus but the primary instrument. A degree-preserving shuffle, a
router-free arm, an architecture-matched untrained floor and a label-shuffled
pipeline control are what separate a mechanism from a coincidence, and each of
them changed a conclusion here.

% ===========================================================================
\section{Limitations}
\label{sec:disc-limits}

\emph{Statistical power.} The study is underpowered for single-component effects
of the size observed. Table~\ref{tab:power} converts each effect into the number
of training seeds required to detect it: one for the full system and for the
ordinal head, ten for typed edges, but 216 for cross-sequence self-supervision
and 3{,}467 for routing. The last is roughly three years of continuous
computation on the hardware used. Those two components are not merely unproven
here; effects of the size observed are not establishable at any plausible scale,
and would carry no clinical meaning if they were.

\emph{External validity is untested.} RQ4 was not executed. The complete system
under institutional transfer, specified as rung E8, was not run. Every result in
this thesis is internal to LumbarDISC, and no claim of generalisation to another
institution, scanner population or reference standard is made.

\emph{Single benchmark, single reference standard.} LumbarDISC labels derive
from a reading process whose inter-reader agreement is moderate for the lateral
compartments. Conclusions about which mechanisms help are conditional on that
reference standard.

\emph{Attribution is not measured on the final system.} The Grad-CAM analysis
covers E0--E4; the graph rungs route through a different batch layout. The
mechanism offered in Section~\ref{sec:disc-localise} is therefore established on
the encoder rather than on E7 in full.

\emph{No clinical outcome is measured.} The study measures agreement with a
retrospective reference standard. It does not measure diagnostic accuracy
against a surgical or clinical ground truth, does not measure effect on
reporting time, and does not support any claim of patient benefit.

\emph{The quality-control reader pass is outstanding.} The ROI inspection
sheets and the checklist of Section~\ref{sec:results-roiqc} have been generated;
the reader adjudication that converts them into the promised
inclusion/exclusion table has not been completed.

% ===========================================================================
\section{Implications and Future Work}
\label{sec:disc-future}

\emph{For anatomy-aware model design.} The most transferable finding is that a
convolutional encoder given a tight, physically-defined crop already performs
most of the localisation that anatomical priors are introduced to supply. Before
adding structure above such an encoder, the marginal room for that structure
should be measured---an attribution floor against an untrained network costs
nothing and would have predicted these results in advance.

\emph{For relational modelling.} A degree-preserving shuffle should be standard
practice wherever a graph's topology is claimed to be meaningful. Without it,
capacity and anatomy are confounded, and this study finds that the capacity
explains the gain.

\emph{For the representation.} The dissociation of
Section~\ref{sec:disc-ordering}---readers find subarticular hardest, the model
finds foraminal hardest---points to the 2.5D representation as a specific,
testable limitation on through-plane structure. A volumetric or genuinely
cross-plane representation is the most direct follow-up this study suggests.

\emph{For transfer.} The infrastructure for RQ4 exists: the cohort is
reconciled, reports are parsed, and the geometric machinery is verified. What
remains is reader adjudication of conflicting reference labels, a correct
de-identification procedure, and---the largest gap---localisation for a cohort
that carries no annotated coordinates. Automated disc localisation is the
enabling step, and it is a separate piece of work rather than an extension of
this one.

% ===========================================================================
\section{Conclusion}
\label{sec:disc-conclusion}

This thesis set out to determine whether explicit anatomical structure improves
automated grading of lumbar degenerative disease. It determined that, at this
data scale and under this reference standard, it largely does not, and it
established why: the convolutional encoder already localises the graded
structure, and the reference standard leaves narrow headroom above it.

The system that was built works, reproducibly and by a margin that survives
correction. The components that make it work are relational typing and a
clinically motivated objective, not anatomical correspondence or
disease-conditioned routing. Two proposed mechanisms are rejected on evidence
gathered specifically to test them, a third is supported in a narrower and more
precise form than proposed, and the study's own controls overturned three
conclusions that would otherwise have been reported as findings.

Chapter~\ref{ch:introduction} framed the aim as \enquote{a determination rather
than a promise of superiority}. That framing was prescient, and this thesis
should be read as having discharged it.

\end{document}
; docmute skips this preamble.
% ===========================================================================

\documentclass[12pt,a4paper]{report}

\usepackage{lmodern}
\usepackage[T1]{fontenc}
\usepackage[utf8]{inputenc}
\usepackage[english]{babel}
\usepackage{csquotes}
\usepackage{microtype}
\usepackage[margin=2.5cm]{geometry}
\usepackage{setspace}
\onehalfspacing
\usepackage{amsmath,amssymb}
\usepackage{booktabs}
\usepackage{array}
\usepackage{graphicx}
\usepackage[hidelinks]{hyperref}

\begin{document}

\chapter{Discussion and Conclusions}
\label{ch:discussion}

% ===========================================================================
\section{What This Study Found}
\label{sec:disc-summary}

The thesis asked whether explicit anatomical structure, imposed on a
multi-sequence grading model, improves automated classification of lumbar
degenerative disease. The answer is that it largely does not, and the more
useful part of the answer is \emph{why}.

The complete system improves on a single-sequence baseline by $+0.0172$ QWK
(95\% CI $[+0.0064, +0.0285]$), reproducibly across seeds and after correction
for multiple comparisons. But the components credited with that improvement in
the original design are not the components responsible for it. Anatomically
aligned self-supervision produced no measurable benefit on any of three axes.
Disease-conditioned routing learned the clinically expected sequence weights
and produced no benefit from having learned them. Relational typing helped;
the anatomical identity of the relations did not.

Two of the three proposed mechanisms are therefore rejected, one is
half-supported in a narrower form than proposed, and the largest single
contributor to final performance is an objective function that
Chapter~\ref{ch:introduction} explicitly declined to claim as original.

This chapter argues that these results are informative rather than merely
disappointing, offers two mechanisms for them, situates them against the
literature of Chapter~\ref{ch:litreview}, and states plainly what the study
cannot conclude.

% ===========================================================================
\section{Why the Anatomical Priors Did Not Pay}
\label{sec:disc-why}

A null result without a mechanism invites the reader to supply their own, and
the most available one---that the implementation was faulty---is the least
interesting and the most damaging. Two mechanisms are supported by evidence
gathered specifically to test them.

\subsection{The Encoder Already Localises}
\label{sec:disc-localise}

The structural priors of RQ1--RQ3 share a premise: that anatomical context is
needed to find and interpret the graded structure, and that supplying it
explicitly will help. The attribution analysis of
Section~\ref{sec:results-attribution} tests that premise directly and finds it
largely unmet.

Every trained configuration concentrates between $2.3$ and $2.9$ times the
chance share of its attention on the annotated lesion. Decisively, E0---a plain
single-sequence convolutional network with no routing, no self-supervision and
no graph---already reaches $0.490$ against an untrained floor of $0.204$ and a
uniform expectation of $0.197$. The localisation the priors were designed to
supply is largely accomplished before any of them is added.

This is consistent with Section~\ref{sec:lr-localize-first}, which
assembles the evidence that localisation precedes classification. The present
result sharpens that literature rather than contradicting it: localisation does
precede classification, and on a benchmark that supplies the localisation as an
annotated coordinate, a convolutional encoder given a tight physical crop
performs the residual localisation unaided. Structure imposed above that
encoder is competing for a job already done.

Section~\ref{sec:intro-conceptual} anticipated this possibility in advance,
stating that \enquote{if local evidence alone is sufficient, more complex
mechanisms should not be retained}. The finding is that local evidence is very
nearly sufficient.

\subsection{The Reference Standard Constrains the Headroom}
\label{sec:disc-ceiling}

Section~\ref{sec:lr-label-ceiling} argues that reader reliability is a hard
constraint on interpretation rather than a peripheral caveat. Lurie et al.\
report inter-reader $\kappa$ of $0.73$ for central canal stenosis, $0.58$ for
foraminal stenosis and $0.49$ for subarticular stenosis
\cite{lurie2008reliability}. A model trained on one reader's labels cannot be
expected to exceed the agreement those readers achieve with each other.

The comparison must be made carefully. Quadratic weighted kappa credits
near-misses that unweighted $\kappa$ penalises fully, so the QWK values reported
here are not directly comparable with Lurie's figures, and the two are measured
on different cohorts under different protocols. The point is not that this
system outperforms radiologists; nothing in this study supports that claim.

The point is about headroom. If the achievable band is bounded by a reference
standard on which experts disagree at $\kappa \approx 0.5$ for the lateral
compartments, then the margin available to any architectural refinement is
narrow, and effects of the size proposed here---$0.005$ QWK---sit far inside the
noise of the labels themselves. Chapter~\ref{ch:litreview} warns that
\enquote{a reported accuracy approaching unity on such a target should prompt
scrutiny of the evaluation protocol rather than admiration}; the corollary is
that failure to improve on such a target by half a percentage point warrants
less alarm than it first appears to.

\subsection{A Test of That Explanation, and Its Limits}
\label{sec:disc-ordering}

Section~\ref{sec:lr-reliability} makes a checkable prediction: agreement is
\enquote{systematically worse for the lateral compartments than for the central
canal---the same ordering that later appears in automated model performance}.
If model difficulty tracks reader difficulty, the label-ceiling explanation
gains support.

\begin{table}[htbp]
\centering
\caption{Per-condition performance against reported inter-reader agreement.
Reader $\kappa$ from Lurie et al.; model QWK averaged over three seeds. The two
metrics are not directly comparable in absolute terms; only the ordering is
being examined.}
\label{tab:disc-ordering}
\begin{tabular}{lcccc}
\toprule
Condition & Reader $\kappa$ & E0 & E2 & E4 \\
\midrule
Central canal        & 0.73 & 0.799 & 0.786 & 0.805 \\
Left foraminal       & 0.58 & 0.659 & 0.623 & 0.644 \\
Right foraminal      & 0.58 & 0.685 & 0.672 & 0.671 \\
Left subarticular    & 0.49 & 0.758 & 0.749 & 0.767 \\
Right subarticular   & 0.49 & 0.712 & 0.731 & 0.736 \\
\midrule
Central $-$ lateral mean & $+0.20$ & $+0.095$ & $+0.092$ & $+0.100$ \\
\bottomrule
\end{tabular}
\end{table}

The coarse prediction holds. Central canal stenosis is the easiest condition for
the model, as it is for readers, and the central-versus-lateral gap runs in the
same direction for both, though smaller in magnitude for the model.

The fine ordering does not hold, and the discrepancy is informative. Readers find
\emph{subarticular} stenosis hardest at $\kappa = 0.49$; the model finds
\emph{foraminal} stenosis hardest, while subarticular is its second-best
condition. Model difficulty is therefore not a simple reflection of label noise.

A representational explanation fits the discrepancy. Foraminal targets are
graded on sagittal T1, where left--right is a \emph{through-plane} axis: the two
sides are different slices rather than different in-plane positions. Subarticular
targets are graded on axial T2, where left--right is in-plane and the lateral
recesses are directly visible. The 2.5D representation used here---three adjacent
slices treated as channels---captures in-plane structure well and through-plane
structure poorly, which is exactly the axis on which foraminal grading depends.
If correct, this locates the residual difficulty in the input representation
rather than in the absence of anatomical priors, and it predicts that a volumetric
or cross-plane-aware representation would help foraminal grading specifically.

That prediction is not tested here. It is offered as the most economical account
of a dissociation the study did not set out to produce, and it is the clearest
direction for subsequent work.

Two cautions apply to this section. Only five conditions exist, two pairs of
which share a reported $\kappa$, so rank correlation between the two orderings
has essentially no statistical power and none is claimed. And Lurie's cohort is
not LumbarDISC; the comparison is indicative rather than a matched analysis.

% ===========================================================================
\section{The Components Individually}
\label{sec:disc-components}

\subsection{Anatomical Self-Supervision (RQ2)}

The pretext task was learnable and was learned: validation InfoNCE reached
$1.5162$ against a chance value of $3.4657$, so the encoders demonstrably
learned to match corresponding anatomy across sequences. The representation
that resulted conferred no measurable downstream benefit on accuracy, on
robustness to a withheld sequence, or on where the model attends.

The natural objection is that the correspondence itself was wrong---that the
pretext task matched misaligned patches and the null is an artefact. That
objection is answered. Section~\ref{sec:results-roiqc} verifies the DICOM
correspondence externally, comparing two independent human annotations through
the transform across 9{,}542 projections, of which 94\% on the pretraining
partition land within one slice of an independently annotated target at the same
level. ACSSL was given correctly aligned data and still conferred nothing.

The result is best read alongside Section~\ref{sec:disc-localise}. Cross-sequence
self-supervision is designed to teach an encoder that one anatomy appears under
several acquisitions. If the supervised task already drives the encoder to
localise that anatomy, the pretext task is teaching something the downstream
objective supplies for free.

\subsection{Disease-Conditioned Routing (RQ3)}

The routing result is the study's clearest illustration of why a mechanism check
requires a control. The router reproduces the clinically expected allocation in
all fifteen runs in which a gate is present, and withholding the
high-weight sequence for a condition is severely damaging---between $0.06$ and
$0.35$ QWK, an order of magnitude larger than any difference in the ablation
ladder. Presented alone, that is a compelling causal story.

It does not survive its control. E1, which has no router at all, exhibits the
same dependence with drops as large or larger. The pattern belongs to the
dataset: each condition is annotated on one particular sequence, so withholding
that sequence removes the only crop centred on the pathology. Any model suffers,
routed or not. Isolating the router's own contribution to sequence reliance
gives $-0.0010 \pm 0.0382$ across seeds, indistinguishable from zero.

Chapter~\ref{ch:methodology} anticipated precisely this trap, stating that
\enquote{the gate weights are model allocations, not causal estimates of
sequence importance}, and requiring the intervention that produced the control.
The methodology's caution was correct and load-bearing.

What robustness the system does exhibit under missing sequences comes from
modality dropout, which reduces reliance on the annotated sequence by
$-0.0588$ on every seed. That is a standard regulariser, disclaimed in
Section~\ref{sec:intro-contrib-supporting} as not original, and it outperforms
the mechanism designed for the purpose.

\subsection{Relational Structure (RQ1)}

RQ1 divides, and the division is the more interesting result. Typed
heterogeneous edges improve on a homogeneous graph by $+0.0123$ QWK on every
seed. Anatomically correct topology does not improve on a degree-preserving
random shuffle: $+0.0051$, two seeds of three, well inside the corrected
threshold. The control preserves node count, edge count and degree, so the two
models are matched in capacity and differ only in \emph{which} targets are
connected to \emph{which}.

The conclusion is therefore narrower and sharper than H3 proposed. What the
graph supplies is a structured parameterisation of interaction between
targets---additional capacity, organised by relation type---rather than the
specific adjacency of levels, compartments and bilateral pairs that motivated
it. Relational typing carries information; anatomical adjacency does not.

Section~\ref{sec:lr-graphs} notes that relational and graph-based modelling of
lumbar targets remains sparsely explored. The present result suggests that where
it is explored, the shuffled-topology control should be mandatory, because a
typed graph will outperform a homogeneous one for reasons that have nothing to
do with the anatomy its edges are named after.

\subsection{Ordinal and Cost-Sensitive Objectives (RQ5)}

RQ5 is the study's most reliable positive result and it contradicts the strongest
prior evidence. Section~\ref{sec:lr-ordinal-losses} reports that Niemeyer et al.\
tested soft kappa loss, ordinal cross-entropy and regression losses against
standard cross-entropy for Pfirrmann grading across 7{,}948 discs with extensive
hyperparameter optimisation, and that none improved classification performance
\cite{niemeyer2021a}.

Here the ordinal and cost-sensitive head improves QWK by $+0.0099$ on every seed,
with the smallest between-seed variance of any comparison in the study
($0.0017$), and raises recall on Severe targets from $59.6\%$ to $65.0\%$
relative to the same architecture under categorical cross-entropy.

The difference from Niemeyer et al.\ is most plausibly the cost matrix rather
than the ordinal structure. Their comparison isolated ordinality; the objective
tested here additionally encodes clinical asymmetry, penalising
Severe$\rightarrow$Normal errors more heavily than the reverse. The gain appears
where that asymmetry acts. This reading is consistent with
Section~\ref{sec:lr-error-asymmetry}, which argues that the clinical cost of an
error depends on its direction while standard objectives optimise a symmetric
one.

It is worth stating plainly that the supporting component outperformed the novel
ones, and that this is a result about where effort is best spent rather than an
embarrassment.

% ===========================================================================
\section{Methodological Implications}
\label{sec:disc-method}

Three analyses in this study produced positive conclusions that their controls
subsequently overturned. All three erred in the direction of the hypothesis
under test. They are reported in Section~\ref{sec:results-threats} and are
summarised here because the pattern is more instructive than any individual
case.

\emph{Per-seed confidence intervals understate variance.} A paired bootstrap
that resamples patients with the trained model held fixed excludes training
stochasticity. On this problem that omission is decisive: one comparison
returned $+0.0270$ with $p = 0.000$ on one seed and $-0.0234$ with $p = 0.000$
on another. Reported per seed, any claim in this study could have been supported
by choosing a seed.

\emph{A mechanism check without a control reads as proof.} The routing
intervention produced 5/5 agreement with effects an order of magnitude above any
ladder difference, and a router-free model reproduced it exactly.

\emph{An instrument can penalise the hypothesis it is measuring.} The first
attribution probe hooked the first encoder unconditionally, which for
multi-sequence configurations reads the sagittal T1 encoder even on targets
graded from axial T2. It reported that no configuration improved on the
baseline; hooking the encoder that actually read each target reverses the sign.

The general lesson is that in a study of this kind the controls are not
supporting apparatus but the primary instrument. A degree-preserving shuffle, a
router-free arm, an architecture-matched untrained floor and a label-shuffled
pipeline control are what separate a mechanism from a coincidence, and each of
them changed a conclusion here.

% ===========================================================================
\section{Limitations}
\label{sec:disc-limits}

\emph{Statistical power.} The study is underpowered for single-component effects
of the size observed. Table~\ref{tab:power} converts each effect into the number
of training seeds required to detect it: one for the full system and for the
ordinal head, ten for typed edges, but 216 for cross-sequence self-supervision
and 3{,}467 for routing. The last is roughly three years of continuous
computation on the hardware used. Those two components are not merely unproven
here; effects of the size observed are not establishable at any plausible scale,
and would carry no clinical meaning if they were.

\emph{External validity is untested.} RQ4 was not executed. The complete system
under institutional transfer, specified as rung E8, was not run. Every result in
this thesis is internal to LumbarDISC, and no claim of generalisation to another
institution, scanner population or reference standard is made.

\emph{Single benchmark, single reference standard.} LumbarDISC labels derive
from a reading process whose inter-reader agreement is moderate for the lateral
compartments. Conclusions about which mechanisms help are conditional on that
reference standard.

\emph{Attribution is not measured on the final system.} The Grad-CAM analysis
covers E0--E4; the graph rungs route through a different batch layout. The
mechanism offered in Section~\ref{sec:disc-localise} is therefore established on
the encoder rather than on E7 in full.

\emph{No clinical outcome is measured.} The study measures agreement with a
retrospective reference standard. It does not measure diagnostic accuracy
against a surgical or clinical ground truth, does not measure effect on
reporting time, and does not support any claim of patient benefit.

\emph{The quality-control reader pass is outstanding.} The ROI inspection
sheets and the checklist of Section~\ref{sec:results-roiqc} have been generated;
the reader adjudication that converts them into the promised
inclusion/exclusion table has not been completed.

% ===========================================================================
\section{Implications and Future Work}
\label{sec:disc-future}

\emph{For anatomy-aware model design.} The most transferable finding is that a
convolutional encoder given a tight, physically-defined crop already performs
most of the localisation that anatomical priors are introduced to supply. Before
adding structure above such an encoder, the marginal room for that structure
should be measured---an attribution floor against an untrained network costs
nothing and would have predicted these results in advance.

\emph{For relational modelling.} A degree-preserving shuffle should be standard
practice wherever a graph's topology is claimed to be meaningful. Without it,
capacity and anatomy are confounded, and this study finds that the capacity
explains the gain.

\emph{For the representation.} The dissociation of
Section~\ref{sec:disc-ordering}---readers find subarticular hardest, the model
finds foraminal hardest---points to the 2.5D representation as a specific,
testable limitation on through-plane structure. A volumetric or genuinely
cross-plane representation is the most direct follow-up this study suggests.

\emph{For transfer.} The infrastructure for RQ4 exists: the cohort is
reconciled, reports are parsed, and the geometric machinery is verified. What
remains is reader adjudication of conflicting reference labels, a correct
de-identification procedure, and---the largest gap---localisation for a cohort
that carries no annotated coordinates. Automated disc localisation is the
enabling step, and it is a separate piece of work rather than an extension of
this one.

% ===========================================================================
\section{Conclusion}
\label{sec:disc-conclusion}

This thesis set out to determine whether explicit anatomical structure improves
automated grading of lumbar degenerative disease. It determined that, at this
data scale and under this reference standard, it largely does not, and it
established why: the convolutional encoder already localises the graded
structure, and the reference standard leaves narrow headroom above it.

The system that was built works, reproducibly and by a margin that survives
correction. The components that make it work are relational typing and a
clinically motivated objective, not anatomical correspondence or
disease-conditioned routing. Two proposed mechanisms are rejected on evidence
gathered specifically to test them, a third is supported in a narrower and more
precise form than proposed, and the study's own controls overturned three
conclusions that would otherwise have been reported as findings.

Chapter~\ref{ch:introduction} framed the aim as \enquote{a determination rather
than a promise of superiority}. That framing was prescient, and this thesis
should be read as having discharged it.

\end{document}
; docmute skips this preamble.
% ===========================================================================

\documentclass[12pt,a4paper]{report}

\usepackage{lmodern}
\usepackage[T1]{fontenc}
\usepackage[utf8]{inputenc}
\usepackage[english]{babel}
\usepackage{csquotes}
\usepackage{microtype}
\usepackage[margin=2.5cm]{geometry}
\usepackage{setspace}
\onehalfspacing
\usepackage{amsmath,amssymb}
\usepackage{booktabs}
\usepackage{array}
\usepackage{graphicx}
\usepackage[hidelinks]{hyperref}

\begin{document}

\chapter{Discussion and Conclusions}
\label{ch:discussion}

% ===========================================================================
\section{What This Study Found}
\label{sec:disc-summary}

The thesis asked whether explicit anatomical structure, imposed on a
multi-sequence grading model, improves automated classification of lumbar
degenerative disease. The answer is that it largely does not, and the more
useful part of the answer is \emph{why}.

The complete system improves on a single-sequence baseline by $+0.0172$ QWK
(95\% CI $[+0.0064, +0.0285]$), reproducibly across seeds and after correction
for multiple comparisons. But the components credited with that improvement in
the original design are not the components responsible for it. Anatomically
aligned self-supervision produced no measurable benefit on any of three axes.
Disease-conditioned routing learned the clinically expected sequence weights
and produced no benefit from having learned them. Relational typing helped;
the anatomical identity of the relations did not.

Two of the three proposed mechanisms are therefore rejected, one is
half-supported in a narrower form than proposed, and the largest single
contributor to final performance is an objective function that
Chapter~\ref{ch:introduction} explicitly declined to claim as original.

This chapter argues that these results are informative rather than merely
disappointing, offers two mechanisms for them, situates them against the
literature of Chapter~\ref{ch:litreview}, and states plainly what the study
cannot conclude.

% ===========================================================================
\section{Why the Anatomical Priors Did Not Pay}
\label{sec:disc-why}

A null result without a mechanism invites the reader to supply their own, and
the most available one---that the implementation was faulty---is the least
interesting and the most damaging. Two mechanisms are supported by evidence
gathered specifically to test them.

\subsection{The Encoder Already Localises}
\label{sec:disc-localise}

The structural priors of RQ1--RQ3 share a premise: that anatomical context is
needed to find and interpret the graded structure, and that supplying it
explicitly will help. The attribution analysis of
Section~\ref{sec:results-attribution} tests that premise directly and finds it
largely unmet.

Every trained configuration concentrates between $2.3$ and $2.9$ times the
chance share of its attention on the annotated lesion. Decisively, E0---a plain
single-sequence convolutional network with no routing, no self-supervision and
no graph---already reaches $0.490$ against an untrained floor of $0.204$ and a
uniform expectation of $0.197$. The localisation the priors were designed to
supply is largely accomplished before any of them is added.

This is consistent with Section~\ref{sec:lr-localize-first}, which
assembles the evidence that localisation precedes classification. The present
result sharpens that literature rather than contradicting it: localisation does
precede classification, and on a benchmark that supplies the localisation as an
annotated coordinate, a convolutional encoder given a tight physical crop
performs the residual localisation unaided. Structure imposed above that
encoder is competing for a job already done.

Section~\ref{sec:intro-conceptual} anticipated this possibility in advance,
stating that \enquote{if local evidence alone is sufficient, more complex
mechanisms should not be retained}. The finding is that local evidence is very
nearly sufficient.

\subsection{The Reference Standard Constrains the Headroom}
\label{sec:disc-ceiling}

Section~\ref{sec:lr-label-ceiling} argues that reader reliability is a hard
constraint on interpretation rather than a peripheral caveat. Lurie et al.\
report inter-reader $\kappa$ of $0.73$ for central canal stenosis, $0.58$ for
foraminal stenosis and $0.49$ for subarticular stenosis
\cite{lurie2008reliability}. A model trained on one reader's labels cannot be
expected to exceed the agreement those readers achieve with each other.

The comparison must be made carefully. Quadratic weighted kappa credits
near-misses that unweighted $\kappa$ penalises fully, so the QWK values reported
here are not directly comparable with Lurie's figures, and the two are measured
on different cohorts under different protocols. The point is not that this
system outperforms radiologists; nothing in this study supports that claim.

The point is about headroom. If the achievable band is bounded by a reference
standard on which experts disagree at $\kappa \approx 0.5$ for the lateral
compartments, then the margin available to any architectural refinement is
narrow, and effects of the size proposed here---$0.005$ QWK---sit far inside the
noise of the labels themselves. Chapter~\ref{ch:litreview} warns that
\enquote{a reported accuracy approaching unity on such a target should prompt
scrutiny of the evaluation protocol rather than admiration}; the corollary is
that failure to improve on such a target by half a percentage point warrants
less alarm than it first appears to.

\subsection{A Test of That Explanation, and Its Limits}
\label{sec:disc-ordering}

Section~\ref{sec:lr-reliability} makes a checkable prediction: agreement is
\enquote{systematically worse for the lateral compartments than for the central
canal---the same ordering that later appears in automated model performance}.
If model difficulty tracks reader difficulty, the label-ceiling explanation
gains support.

\begin{table}[htbp]
\centering
\caption{Per-condition performance against reported inter-reader agreement.
Reader $\kappa$ from Lurie et al.; model QWK averaged over three seeds. The two
metrics are not directly comparable in absolute terms; only the ordering is
being examined.}
\label{tab:disc-ordering}
\begin{tabular}{lcccc}
\toprule
Condition & Reader $\kappa$ & E0 & E2 & E4 \\
\midrule
Central canal        & 0.73 & 0.799 & 0.786 & 0.805 \\
Left foraminal       & 0.58 & 0.659 & 0.623 & 0.644 \\
Right foraminal      & 0.58 & 0.685 & 0.672 & 0.671 \\
Left subarticular    & 0.49 & 0.758 & 0.749 & 0.767 \\
Right subarticular   & 0.49 & 0.712 & 0.731 & 0.736 \\
\midrule
Central $-$ lateral mean & $+0.20$ & $+0.095$ & $+0.092$ & $+0.100$ \\
\bottomrule
\end{tabular}
\end{table}

The coarse prediction holds. Central canal stenosis is the easiest condition for
the model, as it is for readers, and the central-versus-lateral gap runs in the
same direction for both, though smaller in magnitude for the model.

The fine ordering does not hold, and the discrepancy is informative. Readers find
\emph{subarticular} stenosis hardest at $\kappa = 0.49$; the model finds
\emph{foraminal} stenosis hardest, while subarticular is its second-best
condition. Model difficulty is therefore not a simple reflection of label noise.

A representational explanation fits the discrepancy. Foraminal targets are
graded on sagittal T1, where left--right is a \emph{through-plane} axis: the two
sides are different slices rather than different in-plane positions. Subarticular
targets are graded on axial T2, where left--right is in-plane and the lateral
recesses are directly visible. The 2.5D representation used here---three adjacent
slices treated as channels---captures in-plane structure well and through-plane
structure poorly, which is exactly the axis on which foraminal grading depends.
If correct, this locates the residual difficulty in the input representation
rather than in the absence of anatomical priors, and it predicts that a volumetric
or cross-plane-aware representation would help foraminal grading specifically.

That prediction is not tested here. It is offered as the most economical account
of a dissociation the study did not set out to produce, and it is the clearest
direction for subsequent work.

Two cautions apply to this section. Only five conditions exist, two pairs of
which share a reported $\kappa$, so rank correlation between the two orderings
has essentially no statistical power and none is claimed. And Lurie's cohort is
not LumbarDISC; the comparison is indicative rather than a matched analysis.

% ===========================================================================
\section{The Components Individually}
\label{sec:disc-components}

\subsection{Anatomical Self-Supervision (RQ2)}

The pretext task was learnable and was learned: validation InfoNCE reached
$1.5162$ against a chance value of $3.4657$, so the encoders demonstrably
learned to match corresponding anatomy across sequences. The representation
that resulted conferred no measurable downstream benefit on accuracy, on
robustness to a withheld sequence, or on where the model attends.

The natural objection is that the correspondence itself was wrong---that the
pretext task matched misaligned patches and the null is an artefact. That
objection is answered. Section~\ref{sec:results-roiqc} verifies the DICOM
correspondence externally, comparing two independent human annotations through
the transform across 9{,}542 projections, of which 94\% on the pretraining
partition land within one slice of an independently annotated target at the same
level. ACSSL was given correctly aligned data and still conferred nothing.

The result is best read alongside Section~\ref{sec:disc-localise}. Cross-sequence
self-supervision is designed to teach an encoder that one anatomy appears under
several acquisitions. If the supervised task already drives the encoder to
localise that anatomy, the pretext task is teaching something the downstream
objective supplies for free.

\subsection{Disease-Conditioned Routing (RQ3)}

The routing result is the study's clearest illustration of why a mechanism check
requires a control. The router reproduces the clinically expected allocation in
all fifteen runs in which a gate is present, and withholding the
high-weight sequence for a condition is severely damaging---between $0.06$ and
$0.35$ QWK, an order of magnitude larger than any difference in the ablation
ladder. Presented alone, that is a compelling causal story.

It does not survive its control. E1, which has no router at all, exhibits the
same dependence with drops as large or larger. The pattern belongs to the
dataset: each condition is annotated on one particular sequence, so withholding
that sequence removes the only crop centred on the pathology. Any model suffers,
routed or not. Isolating the router's own contribution to sequence reliance
gives $-0.0010 \pm 0.0382$ across seeds, indistinguishable from zero.

Chapter~\ref{ch:methodology} anticipated precisely this trap, stating that
\enquote{the gate weights are model allocations, not causal estimates of
sequence importance}, and requiring the intervention that produced the control.
The methodology's caution was correct and load-bearing.

What robustness the system does exhibit under missing sequences comes from
modality dropout, which reduces reliance on the annotated sequence by
$-0.0588$ on every seed. That is a standard regulariser, disclaimed in
Section~\ref{sec:intro-contrib-supporting} as not original, and it outperforms
the mechanism designed for the purpose.

\subsection{Relational Structure (RQ1)}

RQ1 divides, and the division is the more interesting result. Typed
heterogeneous edges improve on a homogeneous graph by $+0.0123$ QWK on every
seed. Anatomically correct topology does not improve on a degree-preserving
random shuffle: $+0.0051$, two seeds of three, well inside the corrected
threshold. The control preserves node count, edge count and degree, so the two
models are matched in capacity and differ only in \emph{which} targets are
connected to \emph{which}.

The conclusion is therefore narrower and sharper than H3 proposed. What the
graph supplies is a structured parameterisation of interaction between
targets---additional capacity, organised by relation type---rather than the
specific adjacency of levels, compartments and bilateral pairs that motivated
it. Relational typing carries information; anatomical adjacency does not.

Section~\ref{sec:lr-graphs} notes that relational and graph-based modelling of
lumbar targets remains sparsely explored. The present result suggests that where
it is explored, the shuffled-topology control should be mandatory, because a
typed graph will outperform a homogeneous one for reasons that have nothing to
do with the anatomy its edges are named after.

\subsection{Ordinal and Cost-Sensitive Objectives (RQ5)}

RQ5 is the study's most reliable positive result and it contradicts the strongest
prior evidence. Section~\ref{sec:lr-ordinal-losses} reports that Niemeyer et al.\
tested soft kappa loss, ordinal cross-entropy and regression losses against
standard cross-entropy for Pfirrmann grading across 7{,}948 discs with extensive
hyperparameter optimisation, and that none improved classification performance
\cite{niemeyer2021a}.

Here the ordinal and cost-sensitive head improves QWK by $+0.0099$ on every seed,
with the smallest between-seed variance of any comparison in the study
($0.0017$), and raises recall on Severe targets from $59.6\%$ to $65.0\%$
relative to the same architecture under categorical cross-entropy.

The difference from Niemeyer et al.\ is most plausibly the cost matrix rather
than the ordinal structure. Their comparison isolated ordinality; the objective
tested here additionally encodes clinical asymmetry, penalising
Severe$\rightarrow$Normal errors more heavily than the reverse. The gain appears
where that asymmetry acts. This reading is consistent with
Section~\ref{sec:lr-error-asymmetry}, which argues that the clinical cost of an
error depends on its direction while standard objectives optimise a symmetric
one.

It is worth stating plainly that the supporting component outperformed the novel
ones, and that this is a result about where effort is best spent rather than an
embarrassment.

% ===========================================================================
\section{Methodological Implications}
\label{sec:disc-method}

Three analyses in this study produced positive conclusions that their controls
subsequently overturned. All three erred in the direction of the hypothesis
under test. They are reported in Section~\ref{sec:results-threats} and are
summarised here because the pattern is more instructive than any individual
case.

\emph{Per-seed confidence intervals understate variance.} A paired bootstrap
that resamples patients with the trained model held fixed excludes training
stochasticity. On this problem that omission is decisive: one comparison
returned $+0.0270$ with $p = 0.000$ on one seed and $-0.0234$ with $p = 0.000$
on another. Reported per seed, any claim in this study could have been supported
by choosing a seed.

\emph{A mechanism check without a control reads as proof.} The routing
intervention produced 5/5 agreement with effects an order of magnitude above any
ladder difference, and a router-free model reproduced it exactly.

\emph{An instrument can penalise the hypothesis it is measuring.} The first
attribution probe hooked the first encoder unconditionally, which for
multi-sequence configurations reads the sagittal T1 encoder even on targets
graded from axial T2. It reported that no configuration improved on the
baseline; hooking the encoder that actually read each target reverses the sign.

The general lesson is that in a study of this kind the controls are not
supporting apparatus but the primary instrument. A degree-preserving shuffle, a
router-free arm, an architecture-matched untrained floor and a label-shuffled
pipeline control are what separate a mechanism from a coincidence, and each of
them changed a conclusion here.

% ===========================================================================
\section{Limitations}
\label{sec:disc-limits}

\emph{Statistical power.} The study is underpowered for single-component effects
of the size observed. Table~\ref{tab:power} converts each effect into the number
of training seeds required to detect it: one for the full system and for the
ordinal head, ten for typed edges, but 216 for cross-sequence self-supervision
and 3{,}467 for routing. The last is roughly three years of continuous
computation on the hardware used. Those two components are not merely unproven
here; effects of the size observed are not establishable at any plausible scale,
and would carry no clinical meaning if they were.

\emph{External validity is untested.} RQ4 was not executed. The complete system
under institutional transfer, specified as rung E8, was not run. Every result in
this thesis is internal to LumbarDISC, and no claim of generalisation to another
institution, scanner population or reference standard is made.

\emph{Single benchmark, single reference standard.} LumbarDISC labels derive
from a reading process whose inter-reader agreement is moderate for the lateral
compartments. Conclusions about which mechanisms help are conditional on that
reference standard.

\emph{Attribution is not measured on the final system.} The Grad-CAM analysis
covers E0--E4; the graph rungs route through a different batch layout. The
mechanism offered in Section~\ref{sec:disc-localise} is therefore established on
the encoder rather than on E7 in full.

\emph{No clinical outcome is measured.} The study measures agreement with a
retrospective reference standard. It does not measure diagnostic accuracy
against a surgical or clinical ground truth, does not measure effect on
reporting time, and does not support any claim of patient benefit.

\emph{The quality-control reader pass is outstanding.} The ROI inspection
sheets and the checklist of Section~\ref{sec:results-roiqc} have been generated;
the reader adjudication that converts them into the promised
inclusion/exclusion table has not been completed.

% ===========================================================================
\section{Implications and Future Work}
\label{sec:disc-future}

\emph{For anatomy-aware model design.} The most transferable finding is that a
convolutional encoder given a tight, physically-defined crop already performs
most of the localisation that anatomical priors are introduced to supply. Before
adding structure above such an encoder, the marginal room for that structure
should be measured---an attribution floor against an untrained network costs
nothing and would have predicted these results in advance.

\emph{For relational modelling.} A degree-preserving shuffle should be standard
practice wherever a graph's topology is claimed to be meaningful. Without it,
capacity and anatomy are confounded, and this study finds that the capacity
explains the gain.

\emph{For the representation.} The dissociation of
Section~\ref{sec:disc-ordering}---readers find subarticular hardest, the model
finds foraminal hardest---points to the 2.5D representation as a specific,
testable limitation on through-plane structure. A volumetric or genuinely
cross-plane representation is the most direct follow-up this study suggests.

\emph{For transfer.} The infrastructure for RQ4 exists: the cohort is
reconciled, reports are parsed, and the geometric machinery is verified. What
remains is reader adjudication of conflicting reference labels, a correct
de-identification procedure, and---the largest gap---localisation for a cohort
that carries no annotated coordinates. Automated disc localisation is the
enabling step, and it is a separate piece of work rather than an extension of
this one.

% ===========================================================================
\section{Conclusion}
\label{sec:disc-conclusion}

This thesis set out to determine whether explicit anatomical structure improves
automated grading of lumbar degenerative disease. It determined that, at this
data scale and under this reference standard, it largely does not, and it
established why: the convolutional encoder already localises the graded
structure, and the reference standard leaves narrow headroom above it.

The system that was built works, reproducibly and by a margin that survives
correction. The components that make it work are relational typing and a
clinically motivated objective, not anatomical correspondence or
disease-conditioned routing. Two proposed mechanisms are rejected on evidence
gathered specifically to test them, a third is supported in a narrower and more
precise form than proposed, and the study's own controls overturned three
conclusions that would otherwise have been reported as findings.

Chapter~\ref{ch:introduction} framed the aim as \enquote{a determination rather
than a promise of superiority}. That framing was prescient, and this thesis
should be read as having discharged it.

\end{document}
